\chapter{Building Blocks of the Protocol Core}\label{ch:components}

Each component below is a small class or set of functions in \code{include/fuse/proto/}, with its own test file
in \code{tests/proto/}. None of them owns a thread or reads a clock.

\section{UDP socket with batched I/O}\label{sec:udp}

\code{UdpSocket} (\code{udp.hpp}) is a move-only wrapper around one bound UDP socket. Its receive buffers are
large enough for the biggest datagram, so the kernel never truncates a block. Peer addresses are carried in an
opaque, comparable \code{PeerAddr} that may hold an IPv4 or an IPv6 address.

Its two batched operations carry the data path:

\begin{description}[style=nextline,leftmargin=1.2cm]
  \item[\code{send_segmented}] passes the kernel one buffer of back-to-back datagrams in a single
    \code{sendmsg}, tagged with the \code{UDP_SEGMENT} control message (generic segmentation offload,
    Linux~4.18)~\cite{gso}. The kernel or the network card cuts it into ordinary datagrams, so the wire sees
    normal MTU-sized packets.
    \begin{itemize}
      \item If the kernel refuses the call (EIO, ENOTSUP or EINVAL), a flag latches and the caller falls back to
        one \code{sendto} per datagram.
      \item The bulk transfer caps a batch at 64 segments. With a small block size a batch could otherwise
        exceed the kernel's per-call segment limit, and the resulting EINVAL cannot be told apart from ``this
        route has no offload'', so offload would be disabled for the rest of the transfer.
    \end{itemize}
  \item[\code{recv_batch}] receives up to $N$ datagrams in one \code{recvmmsg}~\cite{recvmmsg} with
    \code{MSG_WAITFORONE}, and records the source address of \emph{every} datagram.
    \begin{itemize}
      \item Without \code{MSG_WAITFORONE}, the call would block until all $N$ datagrams had arrived.
      \item The per-datagram source address is what allows path validation to notice a NAT rebinding in the
        middle of a batch.
    \end{itemize}
\end{description}

One system call per datagram limits a modern host to about 213,000 datagrams per second whatever the link
speed; batching amortises that cost.

\code{path_mtu}\vviii{} connects a temporary socket to the destination and reads the kernel's route MTU
(\code{IP_MTU} or \code{IPV6_MTU}). Together with the IP and UDP header overhead (28 bytes for IPv4, 48 for IPv6)
it bounds how far the block size may grow. On macOS the MTU cannot be queried this way, so blocks do not grow
there. Version~8 also opens each sending socket in the destination's address family; earlier versions always
opened IPv4 sockets, which made IPv6 hosts unreachable.

\section{Sender registry and selective acknowledgement}\label{sec:registry}

The \code{SenderRegistry} is the store from which lost blocks are resent. It is a circular array of
$W$ slots (the stream's window size) indexed by $\mathit{seq}\bmod W$. Each slot holds:
\begin{itemize}
  \item the sequence number, send time, byte offset and length;
  \item the block's original flags and a validity bit;
  \item the payload itself, embedded in a 16~KiB array.
\end{itemize}

The registry has three operations:
\begin{itemize}
  \item \code{store} overwrites the slot a sequence number maps to.
  \item \code{lookup} returns the slot only if it is valid \emph{and} still holds exactly that sequence number.
    Otherwise the requested block has already been confirmed or overwritten.
  \item \code{confirm} releases a slot once an acknowledgement covers it.
\end{itemize}
The flags are stored because a retransmission must add the retransmission bit without losing the last-block
bit: a retransmitted last block is still the last block, and the receiver has no other way to learn that the
stream has ended.

\paragraph{Applying an acknowledgement\vviii.} Every sender digests acknowledgements through one function,
\code{apply_ack} (\code{registry.cpp:52}), which works in three steps:
\begin{enumerate}
  \item \emph{Cumulative part.} If the acknowledgement's base exceeds the sender's base, every slot up to
    $\min(\mathit{ack.base}, \mathit{next\_seq})$ is confirmed. The clamp matters: an unauthenticated
    acknowledgement must not confirm blocks that were never sent.
  \item \emph{Selective part.} Every set bit $i$ of the bitmask, enumerated with a count-trailing-zeros
    instruction, confirms block $\mathit{ack.base}+i$ if it is below \code{next_seq}.
  \item \emph{Result.} The function returns the number of payload bytes newly confirmed, and raises
    \code{received_end} to one past the highest block the acknowledgement shows as received.
\end{enumerate}

Before version~8 the bitmask was ignored, so burst losses were repaired by NACKs alone and blocks the receiver
already held were resent. The \code{received_end} value lets the retransmission timer distinguish blocks that
are lost (unacknowledged, below it) from blocks that may still be in flight (above it). This is the same
distinction SACK-based recovery makes in TCP~\cite{rfc2018} and QUIC~\cite{rfc9002}.

\section{Receiver window and loss detection}\label{sec:receiver}

\code{ReceiverStream} tracks one incoming stream. Its state is:
\begin{itemize}
  \item a base sequence number, below which everything has been delivered;
  \item a bitmask of received blocks;
  \item the highest sequence number seen, and the send time to echo;
  \item two timers per window position: when the gap at that position was first seen, and when it was last
    reported.
\end{itemize}
In version~8 the bitmask and the timers are \emph{rings} indexed by $\mathit{seq}\bmod 1024$
(Figure~\ref{fig:ring}). Earlier versions shifted whole arrays every time the base moved, which made each
receive $O(W)$.

\begin{figure}[H]
\centering
\fitwidth{\begin{tikzpicture}[x=6.2mm,y=6.2mm,font=\footnotesize]
  \foreach \i/\c in {0/green!35,1/green!35,2/white,3/green!35,4/green!35,5/white,6/white,7/green!35,8/white,9/white,10/white,11/white,12/white,13/white,14/white,15/white}{
    \draw[fill=\c] (\i,0) rectangle ++(1,1);
  }
  \node at (2.5,0.5) {\textbf{?}}; \node at (5.5,0.5) {\textbf{?}}; \node at (6.5,0.5) {\textbf{?}};
  \draw[arr] (2.5,-1.3) node[below,lbl]{base\\(first gap)} -- (2.5,-0.1);
  \draw[arr] (7.5,2.3) node[above,lbl]{highest received} -- (7.5,1.1);
  \draw[decorate,decoration={brace,amplitude=4pt,mirror}] (2,-0.15) -- (8,-0.15) node[midway,below=5pt,lbl]{gaps eligible for NACK after the reorder delay};
  \draw[decorate,decoration={brace,amplitude=4pt}] (0,1.15) -- (16,1.15) node[midway,above=5pt,lbl] {ring of 1,024 bits, position $=\mathit{seq}\bmod 1024$ (16 shown)};
\end{tikzpicture}}
\caption{The receive window. Green positions have arrived. The base sits on the first gap; positions between
the base and the highest block received are candidate losses.}\label{fig:ring}
\end{figure}

\paragraph{Receiving a block.} \code{on_receive} (\code{receiver.cpp:135}) handles one arriving sequence
number:
\begin{enumerate}
  \item Below the base, it returns \emph{Duplicate}.
  \item A full window or more beyond the base, it returns \emph{OutOfWindow} and counts the overflow.
  \item Already marked received, it returns \emph{Duplicate}.
  \item Otherwise it marks the block received. If the block is a new highest, every position between the
    previous highest and this block has just become a gap, and each gets its first-missing timestamp, exactly
    once.
  \item Finally it slides the base over the contiguous received prefix, clearing each bit and both timers
    exactly once as it passes. This keeps the invariant that positions below the base are always clear, which
    is what makes the ring unambiguous.
\end{enumerate}
Each position is touched a constant number of times over its life, so receiving is $O(1)$ amortised. The
relative bitmask an acknowledgement needs, in which bit $i$ means $\mathit{base}+i$, is assembled from the ring
word by word with two shifts per word.

\paragraph{Reporting losses.} \code{collect_nacks} walks the clear bits between the base and the highest block
received, a word at a time. It reports a gap only if two conditions hold:
\begin{itemize}
  \item the gap has lasted longer than a \emph{reorder delay};
  \item the gap was not reported within a \emph{re-NACK interval}.
\end{itemize}
A stream that is not \code{LOSSLESS} never reports losses. The two conditions exist because a block that is
merely late must not be treated as lost, and re-reporting a gap before its retransmission could possibly arrive
only causes duplicate retransmissions. The bulk-transfer receiver scales both intervals with its own estimate of
the round-trip time (Section~\ref{sec:recovery}).

\section{Congestion control}\label{sec:cc}

Each \code{LOSSLESS} stream has its own \code{CongestionController} (\code{congestion.hpp}). Keeping congestion
state per stream rather than per connection is the deliberate difference from QUIC on which the design rests.
The controller's window is in abstract units: the bulk transfer interprets them as units of the configured
block size\vviii, and the multiplexed transfer as blocks.

\paragraph{RTT estimation.} The controller follows RFC~6298~\cite{rfc6298}. The first sample $R$ sets
$\mathit{SRTT}=R$ and $\mathit{RTTVAR}=R/2$. Each later sample updates
\[
\mathit{RTTVAR}\leftarrow\tfrac34\mathit{RTTVAR}+\tfrac14\lvert \mathit{SRTT}-R\rvert,\qquad
\mathit{SRTT}\leftarrow\tfrac78\mathit{SRTT}+\tfrac18R .
\]

\paragraph{AIMD epochs and slow start.} The controller makes one decision per elapsed SRTT
(Figure~\ref{fig:ccfsm}):
\begin{itemize}
  \item If a loss (a NACK, or the first timeout of a run) was observed during the epoch, the window halves, but
    not below a floor of two, and slow start ends.
  \item In slow start\vviii, a clean epoch doubles the window up to the cap.
  \item In congestion avoidance, a run of clean epochs adds a fixed step. Both transports grow by 16 units after
    one clean epoch.
\end{itemize}
Slow start was added because, without it, the window had to climb from one eighth of the cap a few units per
round trip, which takes many round trips on a long path. It is opt-in, so existing users of the component keep
their behaviour unchanged.

\begin{figure}[H]
\centering
\fitwidth{\begin{tikzpicture}[node distance=2.2cm and 4.2cm, >={Stealth[length=2.2mm]}, every edge/.style={draw,thick}]
  \node[st,minimum width=2.6cm] (ss) {Slow start};
  \node[st,minimum width=2.6cm,right=of ss] (ca) {Congestion\\avoidance};
  \node[lbl,left=1.3cm of ss] (s0) {\code{enable_slow_start}};
  \node[lbl,above=0.9cm of ca] (s1) {default start\\(window $=$ cap/8)};
  \draw[->,thick] (s0) -- (ss);
  \draw[->,thick] (s1) -- (ca);
  \path[->] (ss) edge[loop below] node[lbl,below] {clean epoch: $w\leftarrow\min(2w,\mathit{cap})$} (ss)
            (ss) edge[bend left=12] node[lbl,above] {loss: $w\leftarrow\max(w/2,2)$, or $w=\mathit{cap}$} (ca)
            (ca) edge[loop below] node[lbl,below,align=center] {loss: $w\leftarrow\max(w/2,2)$\\$n$ clean epochs: $w\leftarrow\min(w+\mathit{step},\mathit{cap})$} (ca);
\end{tikzpicture}}
\caption{The congestion controller's states. One decision is made per SRTT epoch.}\label{fig:ccfsm}
\end{figure}

\paragraph{Retransmission timeout\vviii.} The timeout is computed as
\[
\mathit{RTO}=\max\bigl(\mathit{SRTT}+\max(4\,\mathit{RTTVAR},1\,\mathrm{ms}),\;2\,\mathit{SRTT}\bigr)\cdot 2^{b},
\]
clamped to $[5\,\mathrm{ms}, 1\,\mathrm{s}]$. It is 1~s before any sample, and the backoff exponent $b\le6$ is
reset whenever new data is acknowledged. Only the first timeout of a run counts as a loss: an outage of any
length halves the window once, and later timeouts only back off. Each term has a reason:
\begin{itemize}
  \item The variance term replaces a fixed multiple of SRTT, which fired spuriously under jitter.
  \item The $2\,\mathit{SRTT}$ floor exists because NACK-driven repair needs about half a round trip of reorder
    tolerance plus a full round trip; a timer that fires earlier only duplicates that repair. We found this by
    measurement.
  \item Exponential backoff keeps a dead or overloaded path from being probed at a fixed rate.
\end{itemize}

\paragraph{Pacing\vviii.} After sending $b$ bytes with a window of $B$ bytes, the sender waits
$0.8\cdot b\cdot\mathit{SRTT}/B$. This spreads a window over one round trip at 1.25 times the window rate
instead of releasing it as one line-rate burst, which overflows shallow router buffers and turns congestion
control into a loss generator. The factor 1.25 keeps pacing itself from becoming the bottleneck.

\section{Lock-free single-producer, single-consumer ring}

\code{SpscRing<T>} is a bounded ring with exactly one producer and one consumer thread. Its capacity is rounded
up to a power of two so that an index is a bit mask. The head and tail counters increase monotonically and sit
on separate 64-byte cache lines to avoid false sharing. The protocol is:
\begin{enumerate}
  \item The producer fills slots and then publishes the tail with release ordering.
  \item The consumer reads the tail with acquire ordering, reads the slots, and publishes the head with release
    ordering.
  \item The producer reads the head with acquire ordering before reusing a slot.
\end{enumerate}
Each side caches the other's counter and re-reads the shared variable only when the ring appears full or empty,
so in steady state each side touches the other's cache line once per batch rather than once per item.

There is no mutex and no compare-and-swap, so neither side ever waits for the other. A full ring is simply
reported to the producer. The bulk-transfer receiver then drops the datagram and lets the normal recovery
resend it, rather than stall the socket reader behind a slow disk.

\section{Stage 4 primitives: send queue and reassembly}

\code{SendQueue} is a bounded ring of sequence numbers waiting to be sent, and its behaviour depends on the
stream's \code{COALESCE} flag:
\begin{itemize}
  \item \code{COALESCE=1}: when the queue is full, the oldest unsent entry is dropped to admit the newest, and
    the drop is counted. Fresh data supersedes stale data.
  \item \code{COALESCE=0}: a full queue refuses the push, and the producer backs off.
\end{itemize}
A capacity of zero, which used to cause a division by zero, is now clamped to one.

\code{ReassemblyStream} delivers received blocks according to the \code{ORDERED} flag:
\begin{itemize}
  \item \code{ORDERED=1}: blocks are buffered and released in sequence order.
  \item \code{ORDERED=0}: each block is released as soon as it arrives.
\end{itemize}
Both modes write to the same sink at the block's explicit offset, so the completed result is identical and only
the order of delivery differs. Version~8 replaced its unbounded duplicate set with a window-bounded ring.

These two primitives implement the per-stream flags that motivate the whole protocol. Today they are exercised
by the tests and the \code{fuse_bench} harness but not exposed by any public transport
(Chapter~\ref{ch:conclusion}).

\section{Autoscaling worker pool}\label{sec:orch}

\code{WorkerOrchestrator} runs a task supplied by the caller on a pool of threads and resizes the pool the way a
Kubernetes horizontal pod autoscaler does.
\begin{itemize}
  \item Each worker repeatedly calls the task, which reports whether it did useful work, and accumulates its
    busy and total time. An idle call sleeps 200~\textmu s, and the sleep is counted as idle time.
  \item Periodically, \code{tick} computes the mean utilisation since the previous tick and resets the
    counters:
    \begin{itemize}
      \item above 0.70 utilisation, with fewer than the maximum number of workers, it adds a worker;
      \item below 0.30, with more than the minimum, it retires the newest worker, which finishes its current
        call (a graceful drain) and is joined outside the lock;
      \item no two actions may be closer together than a 200~ms stabilisation window.
    \end{itemize}
\end{itemize}
The two separate thresholds and the stabilisation window prevent the pool from oscillating.

This component replaced Stage~6, which pinned each worker to a fixed CPU core. Pinning was measured to be about
6\% slower, with a higher median latency, because pinning the senders while their peers float prevents the
scheduler from placing communicating threads together. Throughput also proved not to be monotonic in the number
of workers: four lanes beat eight on the test host. The right degree of concurrency is something to discover at
run time.

\section{The SETUP handshake}\label{sec:setup}

Before any data flows, the multiplexed transfer agrees on the stream table using a three-way handshake
(Figure~\ref{fig:setup}).
\begin{enumerate}
  \item The initiator sends SETUP\_DATA.
  \item The responder validates it, hashes the body exactly as received, draws a fresh nonce\vviii, and replies
    with SETUP\_HASH\_REPLY(hash, nonce).
  \item If the hash matches its own, the initiator answers with SETUP\_FINACK(hash, nonce), which completes the
    handshake on the responder.
\end{enumerate}
A lost SETUP\_DATA or reply is recovered by the initiator, which retransmits every 50~ms up to five times. A lost
FINACK is recovered by the responder, which re-sends its reply; the initiator answers a duplicate reply with
another FINACK.

\begin{figure}[H]
\centering
\fitwidth{\begin{tikzpicture}[font=\small]
  \node[box] (i) at (0,0) {Initiator (sender)};
  \node[box] (r) at (8,0) {Responder (receiver)};
  \draw[thick,black!40] (0,-0.4) -- (0,-5.2); \draw[thick,black!40] (8,-0.4) -- (8,-5.2);
  \draw[msg] (0,-1) -- node[above,sloped] {SETUP\_DATA(config)} (8,-1.6);
  \node[lbl,right] at (8.1,-1.9) {validate; $h=$FNV-1a(body)\\nonce $=$ random\vviii};
  \draw[msg] (8,-2.6) -- node[above,sloped] {SETUP\_HASH\_REPLY($h$, nonce)} (0,-3.2);
  \node[lbl,left] at (-0.1,-3.4) {$h$ equals hash of\\my config?};
  \draw[msg] (0,-4.0) -- node[above,sloped] {SETUP\_FINACK($h$, nonce)} (8,-4.6);
  \node[lbl,right] at (8.1,-4.8) {$h$ and nonce match:\\complete};
\end{tikzpicture}}
\caption{The SETUP handshake. Retransmission timers on both sides recover any single loss.}\label{fig:setup}
\end{figure}

The responder echoes its \emph{own} hash of what it received rather than a simple ``ok''. A corrupted datagram
therefore produces a mismatched hash, so the handshake does not complete, whereas a corruptible boolean could
not detect that.

The responder nonce\vviii{} closes a replay attack: without it, a recorded SETUP\_DATA and SETUP\_FINACK pair
could be replayed later to complete a new handshake.

The responder also validates the offered configuration and ignores, without ever confirming, a SETUP that
violates any of these rules:
\begin{itemize}
  \item the protocol version must be its own;
  \item stream identifiers must be unique;
  \item block sizes must lie in 1--16,384;
  \item windows must lie in 1--1,024.
\end{itemize}
The multiplexed receiver additionally requires stream identifiers $0,\ldots,n-1$, because it indexes an array
with them. Both endpoints are pure state machines driven by received datagrams and timer calls.

\section{Hashing}

Two hash functions are provided:
\begin{itemize}
  \item \textbf{FNV-1a-64}~\cite{fnv}, a fast non-cryptographic hash. It serves as the SETUP checksum, as the
    file identity, and as the digest fallback in builds without cryptography.
  \item \textbf{SipHash-2-4}~\cite{siphash}, a keyed pseudorandom function. It computes the address-validation
    tokens\vviii, which must be unforgeable without the receiver's secret yet cheap enough to compute for every
    incoming StreamStart, including spoofed ones. SipHash was designed for exactly this kind of short keyed
    input, and it is available even in builds without wolfSSL.
\end{itemize}

\section{Session cryptography}\label{sec:crypto}

\paragraph{Why not one DTLS session per lane.} The obvious way to encrypt a sharded transfer is to give each
lane its own DTLS session. The implementation rejects this for two reasons, documented in
\code{session_crypto.hpp}:
\begin{itemize}
  \item A DTLS session owns record-layer state: sequence numbers, epoch, replay window and cipher state. That
    state cannot be shared across lanes without a mutex, which would serialise exactly the parallelism sharding
    creates.
  \item Writing through DTLS hands the socket to the library one record at a time, which destroys segmentation
    offload.
\end{itemize}
Instead, keys are derived once per session, and every lane encrypts independently with the same session key and
a nonce unique to the pair (lane, sequence number). Encrypted blocks remain fixed-size datagrams and pack into an
offload batch exactly as plaintext does. Fast QUIC implementations have the same shape: per-packet AEAD, followed
by batching of already-encrypted packets.

\paragraph{The lane cipher.} \code{LaneCipher} wraps one AES-256-GCM context. The tag is written directly after
the ciphertext so the datagram stays contiguous. A second initialisation is refused, because re-keying an
existing object is exactly how a (key, nonce) pair could repeat.

\paragraph{Key epochs\vviii.} One AES-GCM key should not protect an unbounded number of records; RFC~9001
applies a confidentiality limit of about $2^{23}$ packets~\cite{rfc9001}. Rather than adding rekeying messages,
each lane's key is derived deterministically from the session key, the lane and the epoch
$\mathit{seq}\gg22$, so both ends change keys at the same sequence number with nothing on the wire. A small
schedule keeps the two most recently used epochs, so that retransmissions from just before a boundary do not
force repeated derivations.

\paragraph{Control-plane authentication\vviii.} A \code{ControlAuth} object per lane and side appends and
verifies the GMAC trailer described in Section~\ref{sec:wirecrypto}, with its own key and a send counter that
starts at one. Replaying an old authentic control message is harmless by construction:
\begin{itemize}
  \item acknowledgement bases only ever move forward;
  \item an abort can only be replayed if one was genuinely sent.
\end{itemize}

\paragraph{Security properties.} The per-session random salt gives \emph{key freshness}: two sessions never
share a key. The pre-shared-key path does \emph{not} give forward secrecy. Anyone who records the traffic and
later obtains the key can decrypt it.

\paragraph{Digest and DTLS.} The digest used to verify resumed files is SHA-256 over the list of
per-64-MiB-segment SHA-256 values plus the file length, which allows the segments to be hashed in parallel.

A DTLS~1.2 wrapper with ECDHE-PSK cipher suites~\cite{rfc6347,rfc4279}, which does give forward secrecy, exists
as a standalone component and is tested, but no transport uses it yet. Version~8 made its debug copy of the wire
traffic opt-in.
